% Imporditud: tools/import_eksperimendid.py
% Allikas: Eesti füüsikaolümpiaadi eksperimendi ülesannete
%   kogu (Rael Kalda, 2023), ülesanne Ü127, lahendus L127.
\setAuthor{EFO žürii}
\setRound{lõppvoor}
\setYear{2001}
\setNumber{G E2}
\setDifficulty{4}
\setTopic{laineoptika}
\setVahendid{CD (või osa sellest), kauge punktvalgusallikas (Päike või laelamp), joonlaud}
\setPunktid{11}
\setAllikas{Eksperimendi kogu Ü127, lahendus lk 99}
\setLahendusseis{toores}

\prob{CD-plaat}
Hinnata CD radade vahekaugust. Vajadusel lugeda sin $\alpha \approx \alpha$, punase valguse lainepikkus $\lambda$p = 600 nm ja violetse valguse lainepikkus $\lambda$v = 400 nm.

\hint

\solu
CD-plaadi pinda võib vaadelda kui peegeldifraktsioonvõret. Suuname sellele kaugest valgusallikast (päike või laelamp) tulevad paralleelsed kiired ja jälgime difrageerunud valgust. Mõõdame ära I järku spektri laiuse CD-l (violetse ja punase serva vahekauguse) ning silma kauguse plaadist. Siis saame arvutada spektri servadele vastavate kiirte vahelise nurga $\alpha$. Edasi tuletame valemi d(sin $\alpha$p $-$ sin $\alpha$v) = m($\lambda$p $- \lambda$v), kus d on otsitav võre samm ja m on spektri järk (antud juhul m = 1). Kasutades lihtsustust sin x $\approx$ x, saame d = m($\lambda$p $- \lambda$v)/$\alpha$. Radade vaheliseks kauguseks saame d $\approx$ 1...2 $\mu$m.
\probend
